\documentclass[10pt,a4paper]{amsart}
\usepackage{amsmath,amssymb,amsthm,graphicx}
\usepackage{verbatim}
\newenvironment{reference}{\comment}{\endcomment}
%\newenvironment{reference}{}{}
\newenvironment{slogan}{\comment}{\endcomment}
\newenvironment{history}{\comment}{\endcomment}
% For commutative diagrams we use Xy-pic
\usepackage[all]{xy}

% We use 2cell for 2-commutative diagrams.
\xyoption{2cell}
% Theorem environments.
%
\theoremstyle{plain}
\newtheorem{theorem}[subsection]{Theorem}
\newtheorem{proposition}[subsection]{Proposition}
\newtheorem{lemma}[subsection]{Lemma}

\theoremstyle{definition}
\newtheorem{definition}[subsection]{Definition}
\newtheorem{example}[subsection]{Example}
\newtheorem{exercise}[subsection]{Exercise}
\newtheorem{situation}[subsection]{Situation}

\theoremstyle{remark}
\newtheorem{remark}[subsection]{Remark}
\newtheorem{remarks}[subsection]{Remarks}

\numberwithin{equation}{subsection}

% Macros
%
\def\lim{\mathop{\mathrm{lim}}\nolimits}
\def\colim{\mathop{\mathrm{colim}}\nolimits}
\def\Spec{\mathop{\mathrm{Spec}}}
\def\Hom{\mathop{\mathrm{Hom}}\nolimits}
\def\Ext{\mathop{\mathrm{Ext}}\nolimits}
\def\SheafHom{\mathop{\mathcal{H}\!\mathit{om}}\nolimits}
\def\SheafExt{\mathop{\mathcal{E}\!\mathit{xt}}\nolimits}
\def\Sch{\mathit{Sch}}
\def\Mor{\mathop{\mathrm{Mor}}\nolimits}
\def\Ob{\mathop{\mathrm{Ob}}\nolimits}
\def\Sh{\mathop{\mathit{Sh}}\nolimits}
\def\NL{\mathop{N\!L}\nolimits}
\def\CH{\mathop{\mathrm{CH}}\nolimits}
\def\proetale{{pro\text{-}\acute{e}tale}}
\def\etale{{\acute{e}tale}}
\def\QCoh{\mathit{QCoh}}
\def\Ker{\mathop{\mathrm{Ker}}}
\def\Im{\mathop{\mathrm{Im}}}
\def\Coker{\mathop{\mathrm{Coker}}}
\def\Coim{\mathop{\mathrm{Coim}}}

% Boxtimes
%
\DeclareMathSymbol{\boxtimes}{\mathbin}{AMSa}{"02}

%
% Macros for moduli stacks/spaces
%
\def\QCohstack{\mathcal{QC}\!\mathit{oh}}
\def\Cohstack{\mathcal{C}\!\mathit{oh}}
\def\Spacesstack{\mathcal{S}\!\mathit{paces}}
\def\Quotfunctor{\mathrm{Quot}}
\def\Hilbfunctor{\mathrm{Hilb}}
\def\Curvesstack{\mathcal{C}\!\mathit{urves}}
\def\Polarizedstack{\mathcal{P}\!\mathit{olarized}}
\def\Complexesstack{\mathcal{C}\!\mathit{omplexes}}
% \Pic is the operator that assigns to X its picard group, usage \Pic(X)
% \Picardstack_{X/B} denotes the Picard stack of X over B
% \Picardfunctor_{X/B} denotes the Picard functor of X over B
\def\Pic{\mathop{\mathrm{Pic}}\nolimits}
\def\Picardstack{\mathcal{P}\!\mathit{ic}}
\def\Picardfunctor{\mathrm{Pic}}
\def\Deformationcategory{\mathcal{D}\!\mathit{ef}}
\usepackage{cleveref}
\crefname{originallemma}{original lemma}{original lemmas}
\Crefname{originallemma}{Original lemma}{Original lemmas}
\makeatletter
\newlabel{coherent-lemma-locally-projective-pushforward}{{02O4}{source}{}{sourcepointer-1}{}}
\newlabel{coherent-lemma-locally-projective-pushforward@cref}{{[originallemma][0][]02O4}{[1][0][]source}{}{}{}}
\makeatother
\crefname{originallemma}{original lemma}{original lemmas}
\Crefname{originallemma}{Original lemma}{Original lemmas}
\makeatletter
\newlabel{lemma-affine-vanishing-higher-direct-images}{{073H}{source}{}{sourcepointer-2}{}}
\newlabel{lemma-affine-vanishing-higher-direct-images@cref}{{[originallemma][0][]073H}{[1][0][]source}{}{}{}}
\makeatother
\crefname{originallemma}{original lemma}{original lemmas}
\Crefname{originallemma}{Original lemma}{Original lemmas}
\makeatletter
\newlabel{lemma-coherent-Noetherian}{{07UB}{source}{}{sourcepointer-3}{}}
\newlabel{lemma-coherent-Noetherian@cref}{{[originallemma][0][]07UB}{[1][0][]source}{}{}{}}
\makeatother
\crefname{originallemma}{original lemma}{original lemmas}
\Crefname{originallemma}{Original lemma}{Original lemmas}
\makeatletter
\newlabel{lemma-coherent-Noetherian-quasi-coherent-sub-quotient}{{07UD}{source}{}{sourcepointer-4}{}}
\newlabel{lemma-coherent-Noetherian-quasi-coherent-sub-quotient@cref}{{[originallemma][0][]07UD}{[1][0][]source}{}{}{}}
\makeatother
\crefname{originallemma}{original lemma}{original lemmas}
\Crefname{originallemma}{Original lemma}{Original lemmas}
\makeatletter
\newlabel{lemma-coherent-abelian-Noetherian}{{07UC}{source}{}{sourcepointer-5}{}}
\newlabel{lemma-coherent-abelian-Noetherian@cref}{{[originallemma][0][]07UC}{[1][0][]source}{}{}{}}
\makeatother
\crefname{originallemma}{original lemma}{original lemmas}
\Crefname{originallemma}{Original lemma}{Original lemmas}
\makeatletter
\newlabel{lemma-kill-by-twisting}{{08AQ}{source}{}{sourcepointer-6}{}}
\newlabel{lemma-kill-by-twisting@cref}{{[originallemma][0][]08AQ}{[1][0][]source}{}{}{}}
\makeatother
\crefname{originallemma}{original lemma}{original lemmas}
\Crefname{originallemma}{Original lemma}{Original lemmas}
\makeatletter
\newlabel{lemma-property-higher-rank-cohomological-variant}{{08AN}{source}{}{sourcepointer-7}{}}
\newlabel{lemma-property-higher-rank-cohomological-variant@cref}{{[originallemma][0][]08AN}{[1][0][]source}{}{}{}}
\makeatother
\crefname{originallemma}{original lemma}{original lemmas}
\Crefname{originallemma}{Original lemma}{Original lemmas}
\makeatletter
\newlabel{lemma-weak-chow}{{089J}{source}{}{sourcepointer-8}{}}
\newlabel{lemma-weak-chow@cref}{{[originallemma][0][]089J}{[1][0][]source}{}{}{}}
\makeatother
\crefname{originallemma}{original lemma}{original lemmas}
\Crefname{originallemma}{Original lemma}{Original lemmas}
\makeatletter
\newlabel{sites-cohomology-lemma-relative-Leray}{{0734}{source}{}{sourcepointer-9}{}}
\newlabel{sites-cohomology-lemma-relative-Leray@cref}{{[originallemma][0][]0734}{[1][0][]source}{}{}{}}
\makeatother
\crefname{originallemma}{original lemma}{original lemmas}
\Crefname{originallemma}{Original lemma}{Original lemmas}
\makeatletter
\newlabel{spaces-morphisms-lemma-closed-immersion-affine}{{07U2}{source}{}{sourcepointer-10}{}}
\newlabel{spaces-morphisms-lemma-closed-immersion-affine@cref}{{[originallemma][0][]07U2}{[1][0][]source}{}{}{}}
\makeatother
\crefname{originallemma}{original lemma}{original lemmas}
\Crefname{originallemma}{Original lemma}{Original lemmas}
\makeatletter
\newlabel{spaces-morphisms-lemma-composition-proper}{{04XY}{source}{}{sourcepointer-11}{}}
\newlabel{spaces-morphisms-lemma-composition-proper@cref}{{[originallemma][0][]04XY}{[1][0][]source}{}{}{}}
\makeatother
\crefname{originallemma}{original lemma}{original lemmas}
\Crefname{originallemma}{Original lemma}{Original lemmas}
\makeatletter
\newlabel{spaces-morphisms-lemma-finite-presentation-noetherian}{{04ZL}{source}{}{sourcepointer-12}{}}
\newlabel{spaces-morphisms-lemma-finite-presentation-noetherian@cref}{{[originallemma][0][]04ZL}{[1][0][]source}{}{}{}}
\makeatother
\crefname{originallemma}{original lemma}{original lemmas}
\Crefname{originallemma}{Original lemma}{Original lemmas}
\makeatletter
\newlabel{spaces-morphisms-lemma-immersion-when-closed}{{04CD}{source}{}{sourcepointer-13}{}}
\newlabel{spaces-morphisms-lemma-immersion-when-closed@cref}{{[originallemma][0][]04CD}{[1][0][]source}{}{}{}}
\makeatother
\crefname{originallemma}{original lemma}{original lemmas}
\Crefname{originallemma}{Original lemma}{Original lemmas}
\makeatletter
\newlabel{spaces-morphisms-lemma-locally-finite-type-locally-noetherian}{{04ZK}{source}{}{sourcepointer-14}{}}
\newlabel{spaces-morphisms-lemma-locally-finite-type-locally-noetherian@cref}{{[originallemma][0][]04ZK}{[1][0][]source}{}{}{}}
\makeatother
\crefname{originallemma}{original lemma}{original lemmas}
\Crefname{originallemma}{Original lemma}{Original lemmas}
\makeatletter
\newlabel{spaces-morphisms-lemma-noetherian-finite-type-finite-presentation}{{06G4}{source}{}{sourcepointer-15}{}}
\newlabel{spaces-morphisms-lemma-noetherian-finite-type-finite-presentation@cref}{{[originallemma][0][]06G4}{[1][0][]source}{}{}{}}
\makeatother
\crefname{originallemma}{original lemma}{original lemmas}
\Crefname{originallemma}{Original lemma}{Original lemmas}
\makeatletter
\newlabel{spaces-morphisms-lemma-semi-diagonal}{{03KO}{source}{}{sourcepointer-16}{}}
\newlabel{spaces-morphisms-lemma-semi-diagonal@cref}{{[originallemma][0][]03KO}{[1][0][]source}{}{}{}}
\makeatother
\crefname{originallemma}{original lemma}{original lemmas}
\Crefname{originallemma}{Original lemma}{Original lemmas}
\makeatletter
\newlabel{spaces-morphisms-lemma-support-finite-type}{{07TZ}{source}{}{sourcepointer-17}{}}
\newlabel{spaces-morphisms-lemma-support-finite-type@cref}{{[originallemma][0][]07TZ}{[1][0][]source}{}{}{}}
\makeatother
\crefname{originallemma}{original lemma}{original lemmas}
\Crefname{originallemma}{Original lemma}{Original lemmas}
\makeatletter
\newlabel{spaces-morphisms-lemma-universally-closed-permanence}{{04NX}{source}{}{sourcepointer-18}{}}
\newlabel{spaces-morphisms-lemma-universally-closed-permanence@cref}{{[originallemma][0][]04NX}{[1][0][]source}{}{}{}}
\makeatother
\crefname{originallemma}{original lemma}{original lemmas}
\Crefname{originallemma}{Original lemma}{Original lemmas}
\makeatletter
\newlabel{spaces-properties-lemma-pushforward-etale-base-change-modules}{{03LX}{source}{}{sourcepointer-19}{}}
\newlabel{spaces-properties-lemma-pushforward-etale-base-change-modules@cref}{{[originallemma][0][]03LX}{[1][0][]source}{}{}{}}
\makeatother
\crefname{originallemma}{original lemma}{original lemmas}
\Crefname{originallemma}{Original lemma}{Original lemmas}
\makeatletter
\newlabel{spaces-properties-lemma-quasi-separated-sober}{{06NJ}{source}{}{sourcepointer-20}{}}
\newlabel{spaces-properties-lemma-quasi-separated-sober@cref}{{[originallemma][0][]06NJ}{[1][0][]source}{}{}{}}
\makeatother
\crefname{originallemma}{original lemma}{original lemmas}
\Crefname{originallemma}{Original lemma}{Original lemmas}
\makeatletter
\newlabel{spaces-properties-lemma-stalk-pullback-quasi-coherent}{{05VQ}{source}{}{sourcepointer-21}{}}
\newlabel{spaces-properties-lemma-stalk-pullback-quasi-coherent@cref}{{[originallemma][0][]05VQ}{[1][0][]source}{}{}{}}
\makeatother
\title{Coherence of the higher direct images of a coherent module along a proper morphism of algebraic spaces}
\author{The Stacks Project Authors}
\date{stacks-project:spaces-cohomology.tex:local-snapshot-20261012; GNU FDL 1.2 or later}
\begin{document}
\maketitle
\noindent\textit{Editorial note.} All delivered mathematical text is the author's own, copied from the pinned Stacks Project chapter file spaces-cohomology.tex; no statement or proof step was rewritten, reordered or removed.\par\medskip
\noindent\textit{Editorial note.} Preamble provenance: the pinned chapter file opens with \textbackslash{}input\{preamble\} and carries no preamble text of its own, so the project's single Stacks-wide preamble staged at sources/stacks-spaces-cohomology/preamble.tex is excerpted verbatim as three spans, lines 8-12 (verbatim package and the project's own reference/slogan/history display environments), lines 14-18 (Xy-pic setup) and lines 154-217 (theorem declarations and mathematical macros). The wrapper class and the amsmath/amssymb/amsthm/graphicx base packages are supplied by the assembler. No macro was invented or redefined.\par\medskip
\noindent\textit{Editorial note.} Same-chapter pointers are located at their \textbackslash{}label lines inside the pinned chapter file; cross-chapter pointers are located in the project's own tag list staged in the snapshot (sources/stacks-chow/tags.txt), at the verbatim lines carrying their canonical tags.\par\medskip
\noindent\textit{Allowed standard prerequisites (editorial).}\begin{itemize}
\item Etale locality of the higher direct images and of the coherence check on the target. Properties of Spaces, label spaces-properties-lemma-pushforward-etale-base-change-modules (invoked at proof line 3850) and Same chapter spaces-cohomology.tex, label lemma-coherent-Noetherian (invoked at proof line 3852). Together they let the proof assume that the target is a Noetherian affine scheme, since the higher direct images commute with etale localization and coherence can be tested etale locally.
\item The three-condition criterion for coherence of the higher direct images over a Noetherian affine base. Same chapter spaces-cohomology.tex, label lemma-property-higher-rank-cohomological-variant (invoked at proof lines 3863 and 3868). It reduces the affine case to three properties of the class of coherent modules: two-out-of-three for short exact sequences, stability under direct sums, and the reduction to a coherent module on a closed reduced irreducible subspace.
\item The weak Chow lemma, producing a proper surjective cover mapping into projective space. Same chapter spaces-cohomology.tex, label lemma-weak-chow (invoked at proof line 3919). Applied to \$Z \textbackslash{}to \textbackslash{}Spec(A)\$ it yields the diagram with \$\textbackslash{}pi : Z' \textbackslash{}to Z\$ proper surjective and an immersion \$i : Z' \textbackslash{}to \textbackslash{}mathbf\{P\}\textasciicircum{}n\_A\$, which becomes a closed immersion by properness.
\item Vanishing of the higher direct images after twisting by a relatively ample invertible sheaf. Same chapter spaces-cohomology.tex, label lemma-kill-by-twisting (invoked at proof line 3944). Applied to \$\textbackslash{}mathcal\{L\}\$ with \$\textbackslash{}pi\$ and with \$g'\$ it gives \$d \textbackslash{}gg 0\$ for which all the higher direct images of \$\textbackslash{}mathcal\{L\}\textasciicircum{}\{\textbackslash{}otimes d\}\$ vanish, which is what makes \$\textbackslash{}mathcal\{G\} = \textbackslash{}pi\_*\textbackslash{}mathcal\{L\}\textasciicircum{}\{\textbackslash{}otimes d\}\$ usable.
\item Coherence of the pushforward of a globally generated sheaf on a closed subscheme of projective space. Cohomology of Schemes, label coherent-lemma-locally-projective-pushforward (invoked at proof line 3966). It shows that \$g\_*\textbackslash{}mathcal\{G\} = (g')\_*\textbackslash{}mathcal\{L\}\textasciicircum{}\{\textbackslash{}otimes d\}\$ is coherent, completing the verification of condition (3).
\item Sandwiching of coherent modules in short exact sequences. Same chapter spaces-cohomology.tex, labels lemma-coherent-abelian-Noetherian and lemma-coherent-Noetherian-quasi-coherent-sub-quotient (invoked at proof lines 3889 and 3890). They show that a module sandwiched between two coherent modules is coherent, which is how condition (1) is verified from the long exact sequence of higher direct images.
\end{itemize}\medskip
\ifdefined\hypertarget\hypertarget{sourcepointer-1}{}\fi
\noindent\textit{Source reference.} Original lemma 02O4: sources/stacks-chow/tags.txt, source lines 3503--3503. Navigation pointer for \textbackslash{}ref\{coherent-lemma-locally-projective-pushforward\} used inside the counted span (line 3965); the argument that consumes it is retained in the span.\par\medskip
\ifdefined\hypertarget\hypertarget{sourcepointer-2}{}\fi
\noindent\textit{Source reference.} Original lemma 073H: sources/stacks-spaces-cohomology/pinned-original-source.tex, source lines 1522--1522. Navigation pointer for \textbackslash{}ref\{lemma-affine-vanishing-higher-direct-images\} used inside the counted span (line 3910); the argument that consumes it is retained in the span.\par\medskip
\ifdefined\hypertarget\hypertarget{sourcepointer-3}{}\fi
\noindent\textit{Source reference.} Original lemma 07UB: sources/stacks-spaces-cohomology/pinned-original-source.tex, source lines 2094--2094. Navigation pointer for \textbackslash{}ref\{lemma-coherent-Noetherian\} used inside the counted span (line 3852); the argument that consumes it is retained in the span.\par\medskip
\ifdefined\hypertarget\hypertarget{sourcepointer-4}{}\fi
\noindent\textit{Source reference.} Original lemma 07UD: sources/stacks-spaces-cohomology/pinned-original-source.tex, source lines 2152--2152. Navigation pointer for \textbackslash{}ref\{lemma-coherent-Noetherian-quasi-coherent-sub-quotient\} used inside the counted span (line 3890); the argument that consumes it is retained in the span.\par\medskip
\ifdefined\hypertarget\hypertarget{sourcepointer-5}{}\fi
\noindent\textit{Source reference.} Original lemma 07UC: sources/stacks-spaces-cohomology/pinned-original-source.tex, source lines 2126--2126. Navigation pointer for \textbackslash{}ref\{lemma-coherent-abelian-Noetherian\} used inside the counted span (line 3889); the argument that consumes it is retained in the span.\par\medskip
\ifdefined\hypertarget\hypertarget{sourcepointer-6}{}\fi
\noindent\textit{Source reference.} Original lemma 08AQ: sources/stacks-spaces-cohomology/pinned-original-source.tex, source lines 3809--3809. Navigation pointer for \textbackslash{}ref\{lemma-kill-by-twisting\} used inside the counted span (line 3944); the argument that consumes it is retained in the span.\par\medskip
\ifdefined\hypertarget\hypertarget{sourcepointer-7}{}\fi
\noindent\textit{Source reference.} Original lemma 08AN: sources/stacks-spaces-cohomology/pinned-original-source.tex, source lines 2779--2779. Navigation pointer for \textbackslash{}ref\{lemma-property-higher-rank-cohomological-variant\} used inside the counted span (line 3863); the argument that consumes it is retained in the span.\par\medskip
\ifdefined\hypertarget\hypertarget{sourcepointer-8}{}\fi
\noindent\textit{Source reference.} Original lemma 089J: sources/stacks-spaces-cohomology/pinned-original-source.tex, source lines 3498--3498. Navigation pointer for \textbackslash{}ref\{lemma-weak-chow\} used inside the counted span (line 3919); the argument that consumes it is retained in the span.\par\medskip
\ifdefined\hypertarget\hypertarget{sourcepointer-9}{}\fi
\noindent\textit{Source reference.} Original lemma 0734: sources/stacks-chow/tags.txt, source lines 9048--9048. Navigation pointer for \textbackslash{}ref\{sites-cohomology-lemma-relative-Leray\} used inside the counted span (line 3908); the argument that consumes it is retained in the span.\par\medskip
\ifdefined\hypertarget\hypertarget{sourcepointer-10}{}\fi
\noindent\textit{Source reference.} Original lemma 07U2: sources/stacks-chow/tags.txt, source lines 9957--9957. Navigation pointer for \textbackslash{}ref\{spaces-morphisms-lemma-closed-immersion-affine\} used inside the counted span (line 3913); the argument that consumes it is retained in the span.\par\medskip
\ifdefined\hypertarget\hypertarget{sourcepointer-11}{}\fi
\noindent\textit{Source reference.} Original lemma 04XY: sources/stacks-chow/tags.txt, source lines 6411--6411. Navigation pointer for \textbackslash{}ref\{spaces-morphisms-lemma-composition-proper\} used inside the counted span (line 3929); the argument that consumes it is retained in the span.\par\medskip
\ifdefined\hypertarget\hypertarget{sourcepointer-12}{}\fi
\noindent\textit{Source reference.} Original lemma 04ZL: sources/stacks-chow/tags.txt, source lines 6469--6469. Navigation pointer for \textbackslash{}ref\{spaces-morphisms-lemma-finite-presentation-noetherian\} used inside the counted span (line 3862); the argument that consumes it is retained in the span.\par\medskip
\ifdefined\hypertarget\hypertarget{sourcepointer-13}{}\fi
\noindent\textit{Source reference.} Original lemma 04CD: sources/stacks-chow/tags.txt, source lines 5687--5687. Navigation pointer for \textbackslash{}ref\{spaces-morphisms-lemma-immersion-when-closed\} used inside the counted span (line 3933); the argument that consumes it is retained in the span.\par\medskip
\ifdefined\hypertarget\hypertarget{sourcepointer-14}{}\fi
\noindent\textit{Source reference.} Original lemma 04ZK: sources/stacks-chow/tags.txt, source lines 6468--6468. Navigation pointer for \textbackslash{}ref\{spaces-morphisms-lemma-locally-finite-type-locally-noetherian\} used inside the counted span (line 3846); the argument that consumes it is retained in the span.\par\medskip
\ifdefined\hypertarget\hypertarget{sourcepointer-15}{}\fi
\noindent\textit{Source reference.} Original lemma 06G4: sources/stacks-chow/tags.txt, source lines 8273--8273. Navigation pointer for \textbackslash{}ref\{spaces-morphisms-lemma-noetherian-finite-type-finite-presentation\} used inside the counted span (line 3859); the argument that consumes it is retained in the span.\par\medskip
\ifdefined\hypertarget\hypertarget{sourcepointer-16}{}\fi
\noindent\textit{Source reference.} Original lemma 03KO: sources/stacks-chow/tags.txt, source lines 4684--4684. Navigation pointer for \textbackslash{}ref\{spaces-morphisms-lemma-semi-diagonal\} used inside the counted span (line 3937); the argument that consumes it is retained in the span.\par\medskip
\ifdefined\hypertarget\hypertarget{sourcepointer-17}{}\fi
\noindent\textit{Source reference.} Original lemma 07TZ: sources/stacks-chow/tags.txt, source lines 9954--9954. Navigation pointer for \textbackslash{}ref\{spaces-morphisms-lemma-support-finite-type\} used inside the counted span (line 4006); the argument that consumes it is retained in the span.\par\medskip
\ifdefined\hypertarget\hypertarget{sourcepointer-18}{}\fi
\noindent\textit{Source reference.} Original lemma 04NX: sources/stacks-chow/tags.txt, source lines 6095--6095. Navigation pointer for \textbackslash{}ref\{spaces-morphisms-lemma-universally-closed-permanence\} used inside the counted span (line 3932); the argument that consumes it is retained in the span.\par\medskip
\ifdefined\hypertarget\hypertarget{sourcepointer-19}{}\fi
\noindent\textit{Source reference.} Original lemma 03LX: sources/stacks-chow/tags.txt, source lines 4729--4729. Navigation pointer for \textbackslash{}ref\{spaces-properties-lemma-pushforward-etale-base-change-modules\} used inside the counted span (line 3850); the argument that consumes it is retained in the span.\par\medskip
\ifdefined\hypertarget\hypertarget{sourcepointer-20}{}\fi
\noindent\textit{Source reference.} Original lemma 06NJ: sources/stacks-chow/tags.txt, source lines 8538--8538. Navigation pointer for \textbackslash{}ref\{spaces-properties-lemma-quasi-separated-sober\} used inside the counted span (line 3981); the argument that consumes it is retained in the span.\par\medskip
\ifdefined\hypertarget\hypertarget{sourcepointer-21}{}\fi
\noindent\textit{Source reference.} Original lemma 05VQ: sources/stacks-chow/tags.txt, source lines 7559--7559. Navigation pointer for \textbackslash{}ref\{spaces-properties-lemma-stalk-pullback-quasi-coherent\} used inside the counted span (line 3995); the argument that consumes it is retained in the span.\par\medskip
\begin{lemma}
\label{lemma-proper-pushforward-coherent}
Let $S$ be a scheme. Let $f : X \to Y$ be a proper morphism
of algebraic spaces over $S$ with $Y$ locally Noetherian.
Let $\mathcal{F}$ be a coherent $\mathcal{O}_X$-module.
Then $R^if_*\mathcal{F}$ is a coherent $\mathcal{O}_Y$-module
for all $i \geq 0$.
\end{lemma}

\begin{proof}
We first remark that $X$ is a locally Noetherian algebraic space
by Morphisms of Spaces, Lemma
\ref{spaces-morphisms-lemma-locally-finite-type-locally-noetherian}.
Hence the statement of the lemma makes sense. Moreover, computing
$R^if_*\mathcal{F}$ commutes with \'etale localization on $Y$
(Properties of Spaces, Lemma
\ref{spaces-properties-lemma-pushforward-etale-base-change-modules})
and checking whether $R^if_*\mathcal{F}$ coherent can be done
\'etale locally on $Y$ (Lemma \ref{lemma-coherent-Noetherian}).
Hence we may assume that $Y = \Spec(A)$ is a Noetherian affine scheme.

\medskip\noindent
Assume $Y = \Spec(A)$ is an affine scheme. Note that $f$ is locally
of finite presentation
(Morphisms of Spaces, Lemma
\ref{spaces-morphisms-lemma-noetherian-finite-type-finite-presentation}).
Thus it is of finite presentation, hence $X$ is Noetherian
(Morphisms of Spaces, Lemma
\ref{spaces-morphisms-lemma-finite-presentation-noetherian}).
Thus Lemma \ref{lemma-property-higher-rank-cohomological-variant}
applies to the category of coherent modules of $X$.
For a coherent sheaf $\mathcal{F}$ on $X$ we say $\mathcal{P}$ holds
if and only if $R^if_*\mathcal{F}$ is a coherent module on $\Spec(A)$.
We will show that conditions (1), (2), and (3) of
Lemma \ref{lemma-property-higher-rank-cohomological-variant} hold
for this property thereby finishing the proof of the lemma.

\medskip\noindent
Verification of condition (1). Let
$$
0 \to \mathcal{F}_1 \to \mathcal{F}_2 \to \mathcal{F}_3 \to 0
$$
be a short exact sequence of coherent sheaves on $X$.
Consider the long exact sequence of higher direct images
$$
R^{p - 1}f_*\mathcal{F}_3 \to
R^pf_*\mathcal{F}_1 \to
R^pf_*\mathcal{F}_2 \to
R^pf_*\mathcal{F}_3 \to
R^{p + 1}f_*\mathcal{F}_1
$$
Then it is clear that if 2-out-of-3 of the sheaves $\mathcal{F}_i$
have property $\mathcal{P}$, then the higher direct images of the
third are sandwiched in this exact complex between two coherent
sheaves. Hence these higher direct images are also coherent by
Lemmas \ref{lemma-coherent-abelian-Noetherian} and
\ref{lemma-coherent-Noetherian-quasi-coherent-sub-quotient}.
Hence property $\mathcal{P}$ holds for the third as well.

\medskip\noindent
Verification of condition (2). This follows immediately from the fact
that $R^if_*(\mathcal{F}_1 \oplus \mathcal{F}_2) =
R^if_*\mathcal{F}_1 \oplus R^if_*\mathcal{F}_2$ and that a summand
of a coherent module is coherent (see lemmas cited above).

\medskip\noindent
Verification of condition (3). Let $i : Z \to X$ be a closed immersion
with $Z$ reduced and $|Z|$ irreducible. Set $g = f \circ i : Z \to \Spec(A)$.
Let $\mathcal{G}$ be a coherent module on $Z$ whose scheme theoretic support
is equal to $Z$ such that $R^pg_*\mathcal{G}$ is coherent for all $p$.
Then $\mathcal{F} = i_*\mathcal{G}$ is a coherent module on
$X$ whose scheme theoretic support is $Z$ such that
$R^pf_*\mathcal{F} = R^pg_*\mathcal{G}$. To see this use
the Leray spectral sequence
(Cohomology on Sites, Lemma \ref{sites-cohomology-lemma-relative-Leray})
and the fact that $R^qi_*\mathcal{G} = 0$ for $q > 0$ by
Lemma \ref{lemma-affine-vanishing-higher-direct-images}
and the fact that a closed immersion is affine.
(Morphisms of Spaces, Lemma
\ref{spaces-morphisms-lemma-closed-immersion-affine}).
Thus we reduce to finding a coherent sheaf $\mathcal{G}$ on $Z$
with support equal to $Z$ such that $R^pg_*\mathcal{G}$ is coherent
for all $p$.

\medskip\noindent
We apply Lemma \ref{lemma-weak-chow} to the morphism $Z \to \Spec(A)$.
Thus we get a diagram
$$
\xymatrix{
Z \ar[rd]_g & Z' \ar[d]^-{g'} \ar[l]^\pi \ar[r]_i & \mathbf{P}^n_A \ar[dl] \\
& \Spec(A) &
}
$$
with $\pi : Z' \to Z$ proper surjective and $i$ an immersion.
Since $Z \to \Spec(A)$ is proper we conclude that $g'$ is proper
(Morphisms of Spaces, Lemma \ref{spaces-morphisms-lemma-composition-proper}).
Hence $i$ is a closed immersion
(Morphisms of Spaces, Lemmas
\ref{spaces-morphisms-lemma-universally-closed-permanence} and
\ref{spaces-morphisms-lemma-immersion-when-closed}).
It follows that the morphism
$i' = (i, \pi) : \mathbf{P}^n_A \times_{\Spec(A)} Z' = \mathbf{P}^n_Z$ is
a closed immersion
(Morphisms of Spaces, Lemma \ref{spaces-morphisms-lemma-semi-diagonal}).
Set
$$
\mathcal{L} =
i^*\mathcal{O}_{\mathbf{P}^n_A}(1) =
(i')^*\mathcal{O}_{\mathbf{P}^n_Z}(1)
$$
We may apply Lemma \ref{lemma-kill-by-twisting}
to $\mathcal{L}$ and $\pi$ as well as $\mathcal{L}$ and $g'$.
Hence for all $d \gg 0$ we have
$R^p\pi_*\mathcal{L}^{\otimes d} = 0$ for all $p > 0$ and
$R^p(g')_*\mathcal{L}^{\otimes d} = 0$ for all $p > 0$.
Set $\mathcal{G} = \pi_*\mathcal{L}^{\otimes d}$.
By the Leray spectral sequence
(Cohomology on Sites, Lemma \ref{sites-cohomology-lemma-relative-Leray})
we have
$$
E_2^{p, q} = R^pg_* R^q\pi_*\mathcal{L}^{\otimes d}
\Rightarrow
R^{p + q}(g')_*\mathcal{L}^{\otimes d}
$$
and by choice of $d$ the only nonzero terms in $E_2^{p, q}$ are
those with $q = 0$ and the only nonzero terms of
$R^{p + q}(g')_*\mathcal{L}^{\otimes d}$ are those with $p = q = 0$.
This implies that $R^pg_*\mathcal{G} = 0$ for $p > 0$ and
that $g_*\mathcal{G} = (g')_*\mathcal{L}^{\otimes d}$.
Applying
Cohomology of Schemes, Lemma
\ref{coherent-lemma-locally-projective-pushforward}
we see that $g_*\mathcal{G} = (g')_*\mathcal{L}^{\otimes d}$ is
coherent.

\medskip\noindent
We still have to check that the support of $\mathcal{G}$ is $Z$.
This follows from the fact that $\mathcal{L}^{\otimes d}$ has
lots of global sections. We spell it out here.
Note that $\mathcal{L}^{\otimes d}$ is globally generated for all $d \geq 0$
because the same is true for $\mathcal{O}_{\mathbf{P}^n}(d)$.
Pick a point $z \in Z'$ mapping to the generic point $\xi$ of $Z$
which we can do as $\pi$ is surjective.
(Observe that $Z$ does indeed have a generic point as $|Z|$ is irreducible
and $Z$ is Noetherian, hence quasi-separated, hence $|Z|$ is a sober
topological space by
Properties of Spaces, Lemma
\ref{spaces-properties-lemma-quasi-separated-sober}.)
Pick $s \in \Gamma(Z', \mathcal{L}^{\otimes d})$ which does not vanish at $z$.
Since $\Gamma(Z, \mathcal{G}) = \Gamma(Z', \mathcal{L}^{\otimes d})$
we may think of $s$ as a global section of $\mathcal{G}$.
Choose a geometric point $\overline{z}$ of $Z'$ lying over $z$
and denote $\overline{\xi} = g' \circ \overline{z}$
the corresponding geometric point of $Z$. The adjunction map
$$
(g')^*\mathcal{G} =
(g')^*g'_*\mathcal{L}^{\otimes d} \longrightarrow \mathcal{L}^{\otimes d}
$$
induces a map of stalks
$\mathcal{G}_{\overline{\xi}} \to \mathcal{L}_{\overline{z}}$, see
Properties of Spaces, Lemma
\ref{spaces-properties-lemma-stalk-pullback-quasi-coherent}.
Moreover the adjunction map sends the pullback of $s$ (viewed as a section
of $\mathcal{G}$) to $s$ (viewed as a section of $\mathcal{L}^{\otimes d}$).
Thus the image of $s$ in the vector space which is the source of the arrow
$$
\mathcal{G}_{\overline{\xi}} \otimes \kappa(\overline{\xi})
\longrightarrow
\mathcal{L}^{\otimes d}_{\overline{z}} \otimes \kappa(\overline{z})
$$
isn't zero since by choice of $s$ the image in the target of the arrow
is nonzero. Hence $\xi$ is in the support of $\mathcal{G}$
(Morphisms of Spaces, Lemma \ref{spaces-morphisms-lemma-support-finite-type}).
Since $|Z|$ is irreducible and $Z$ is reduced we conclude that
the scheme theoretic support of $\mathcal{G}$ is all of $Z$ as desired.
\end{proof}

\end{document}
